% Mapped access generalizes to endpoint-compatible I/O pools; network details match the TCP prototype.
\section{Mapped Mode: Segregated I/O Memory}
\label{sec:mapped}

Mapped mode gives a trusted process direct access to buffers in a dedicated I/O area.
A store can lend the same handle and bounds as in named mode, plus offsets into the
area. A network driver can fill eligible pages with TCP or UDP data. A block endpoint
can fill or consume pages laid out for direct I/O. The area and its trust boundary
depend on the endpoint.

\subsection{The memory boundary}
\label{sec:mapped:segregation}

The kernel allocates accounted, zeroed pages for the area. Socket and device state,
I/O descriptors, buffer bookkeeping and object metadata remain in private memory.
The kernel checks object bounds and endpoint arguments against that private state.

In a network area, packet headers can arrive in the same buffers as payload. The
kernel uses private copies for parsing. The virtio-net integration copies protocol
headers into private packet storage and snapshots device control fields, including the
buffer count. Original header bytes can still be visible in the mapping. For block I/O
the kernel similarly keeps the device window, request descriptors and scatter/gather
layout private while the device accesses the data pages.

\subsection{Setup and trust}
\label{sec:mapped:trust}

For networking, a process with \capnetadmin registers an instance on the device. The
kernel reserves its locked-memory allowance and returns an offset for a read/write
mapping. Driver and provider code must supply these pages to the receive queues and
keep the kernel's parsing copies private. Software integration replaces the need
for hardware header splitting and per-application queue steering.

The mapping exposes the area as a whole. In the network prototype it supplies every
receive queue of the device, so a mapper can see or change other flows' data,
including bytes retained by an unbound store. Registration and instance import require
\capnetadmin. Passing the registered ring descriptor can delegate access to the area.
For a block area, trust covers the contents of that area, and an application still
needs authority over the separate device window to issue a read or write.

References keep memory alive. They do not freeze its contents. A mapper can finish
edits before handing an object to a send or write, then keep it stable while that
operation or its endpoint still uses the bytes. A request's completion grants reuse
only when the endpoint has finished using the data. TCP can keep using a page after
the send completion.

\subsection{Raw buffers and retained objects}
\label{sec:mapped:objects}

The network area supports both raw buffer leases and retained objects. A raw receiver
gets offsets and lengths, then returns its receive references through the refill
ring. A raw relay waits for the reuse notification after a zero-copy send before
returning the corresponding receive lease.

With a bound store, an application gets a handle and can query the offsets of a
completed mapped object. The offsets are valid while the application keeps the object
and are not independent references. The kernel returns receive references when the
object and its outstanding users are done. Successful mapped capture puts the
entire object in the area. It retains eligible pages by reference and copies
other readable network data into the area. The same store can also hold unmapped
objects, for example from an ordinary \code{CREATE}.

\subsection{When storage becomes reusable}
\label{sec:mapped:recycle}

A page may be held by the area, by object segments and by I/O outside the store.
Network receive pools track raw consumers and object references separately from
ordinary page references. A page can return to the device when its reference count
returns to the area's and mappings' baseline and all receive-pool users have released
it. Dropping a handle alone cannot make it available for another packet.

After a handle is removed and admitted sends finish queueing, mapped network objects
can move to a waiting list. A background reclaimer passes over objects whose pages
are still busy. Once every page of an object has returned to its baseline reference
count, the reclaimer releases the object's receive-pool references and metadata.
The network page pool exchanges batches using a per-CPU pair with bottom halves
disabled. The store uses separate task and soft-interrupt pairs
(\S\ref{sec:lending:objects}).

Block operations must similarly retain their buffers through completion and wait
for outstanding consumers before recycling a page. A block read publishes a
completed object before the application is entitled to use its offsets, even
though the mapping itself persists. The area bounds the buffers available to
endpoints independently of the store's object limit.
